\chapter{Introduction}

\section{Problem Context}
This thesis studies next-activity prediction in business process event logs using predictive process mining and large language models. The practical goal is to estimate the most likely next step of an ongoing case from its observed prefix.

\section{Motivation}
Traditional predictive process monitoring methods rely on handcrafted representations or specialized sequence models. In contrast, Retrieval-Augmented Generation offers a way to support an LLM with semantically similar historical prefixes and traces, potentially improving prediction quality and interpretability.

\section{Objectives}
The main objectives of this thesis are to design, implement, and evaluate an LLM-based next-activity prediction pipeline with retrieval, and to study how retrieval, prompting, and experimental configuration choices affect performance.

\section{Contributions}
This section will summarize the concrete technical and experimental contributions of the thesis after the final evaluation setup is fixed.

\section{Thesis Structure}
The remainder of the thesis is organized into related work, methodology, experiments, results, and conclusion.

\blankpage
